\documentclass[11pt]{article}
\usepackage[T1]{fontenc}
\usepackage[utf8]{inputenc}
\usepackage{lmodern}
\usepackage[margin=1in]{geometry}
\usepackage{amsmath,amssymb,amsthm,mathtools,bm}
\usepackage{booktabs,array,tabularx}
\usepackage{graphicx}
\usepackage[font=small,labelfont=bf]{caption}
\usepackage{setspace}
\usepackage{microtype}
\usepackage[hidelinks]{hyperref}
\urlstyle{same}
\usepackage[section]{placeins}
\onehalfspacing
\setlength{\emergencystretch}{2em}
\DeclareMathOperator{\tr}{tr}
\DeclareMathOperator{\diag}{diag}
\DeclareMathOperator{\KL}{KL}
\DeclareMathOperator{\TC}{TC}
\DeclareMathOperator{\Var}{Var}
\DeclareMathOperator{\Cov}{Cov}
\newcolumntype{L}{>{\raggedright\arraybackslash}X}
\graphicspath{{figures/}}
\newtheorem{theorem}{Theorem}
\newtheorem{proposition}{Proposition}
\newtheorem{lemma}{Lemma}
\newtheorem{corollary}{Corollary}
\numberwithin{equation}{section}
\hypersetup{
  pdftitle={Fitting One Factor to Two Traits: Population Solutions and Fit Indices},
  pdfauthor={},
  pdfsubject={Manuscript submitted to Psychometrika}
}
